%%%%%%%%%%%%%%%%%%%%%%%%%%%%%%%%%%%%%%%%%
% Cover Letter - 1 page (condensed)
%%%%%%%%%%%%%%%%%%%%%%%%%%%%%%%%%%%%%%%%%

\documentclass{article}

\usepackage{charter}

\usepackage[
	a4paper,
	top=0.4in,
	bottom=0.4in,
	left=1in,
	right=1in,
]{geometry}

\setlength{\parindent}{0pt}
\setlength{\parskip}{0.4em}

\usepackage{graphicx}
\usepackage[colorlinks=true, linkcolor=magenta, citecolor=magenta, urlcolor=magenta]{hyperref}
\usepackage{fancyhdr}

\fancypagestyle{firstpage}{%
	\fancyhf{}
	\renewcommand{\headrulewidth}{0pt}
	\renewcommand{\footrulewidth}{1pt}
}

\AtBeginDocument{\thispagestyle{firstpage}}

\begin{document}

\vspace{-1em}

\rule{\linewidth}{1pt}

\bigskip

\hfill
\begin{tabular}{l @{}}
	\today \bigskip\\
Mahmood Amintoosi\\
Computer Science Division\\
Department of Applied Mathematics \&\\
\href{https://gta-lab.github.io/}{Graph Theory and its Applications Lab} \\
Faculty of Mathematical Sciences,\\ Ferdowsi University of Mashhad,\\
IRAN,\\
\url{https://mamintoosi.github.io/}
\end{tabular}

\bigskip

Dear Editor,

I am submitting my original paper, ``Adaptive-OCGCN: Ambiguity-Aware Overlap Assignment for Scalable Graph Neural Networks,'' for publication in \textit{Progress in Artificial Intelligence}.

Graph Convolutional Networks (GCNs) face scalability challenges on large graphs. Cluster-based training partitions graphs into subgraphs, but hard partitioning discards cross-cluster information. Overlapping Cluster-GCN (OCGCN) allows boundary nodes to participate in multiple clusters using a fixed global threshold, which ignores node-level variation in community certainty. We propose Adaptive-OCGCN, an ambiguity-aware framework that replaces global thresholding with node-specific adaptive decisions. Three strategies are introduced: Entropy-Adaptive, Margin-Adaptive, and Hybrid-Adaptive WMC.

Contributions: (1)~identifying the limitation of fixed global overlap criteria; (2)~proposing Adaptive-OCGCN with node-specific adaptive strategies; (3)~extensive experiments on six benchmarks (Cora, CiteSeer, PubMed, ACM, DBLP, IMDB) showing that adaptive assignment matches tuned fixed-threshold OCGCN while removing dataset-specific WMC tuning; and (4)~ambiguity analysis confirming that high-entropy nodes are 41--95\% overlap candidates.

This work relates to recent graph learning research in this journal. Bo and Yu~\cite{Bo2025} proposed multi-classifier fusion for GCN-based node classification. Khaledian and Falahati~\cite{Khaledian2026} employed DDI-aware cross-attention in heterogeneous graph learning. Our paper complements these by introducing node-level adaptive thresholding for overlap assignment in scalable GNN training.

The paper is well-suited for \textit{Progress in Artificial Intelligence} as it addresses scalable graph representation learning with comprehensive experiments, ablation studies, statistical tests, and ambiguity analysis, aligning with the journal's interest in computational intelligence and machine learning.

The source code is publicly available at \url{https://github.com/mamintoosi/Adaptive-OCGCN}. The datasets are public benchmarks (Cora, CiteSeer, PubMed, ACM, DBLP, IMDB) and are not included in the repository due to their size (~2.5\,GB total), but they can be downloaded from their original sources.

\bigskip

Sincerely,

Mahmood Amintoosi

\begin{thebibliography}{2}
\bibitem{Bo2025} Bo, X., Yu, Y. (2025). Improving node classification using GCN and fused classifiers. \textit{Prog. Artif. Intell.}, 14(3), 355--370.
\bibitem{Khaledian2026} Khaledian, N., Falahati, S. (2026). Safe drug recommendation via heterogeneous graph learning. \textit{Prog. Artif. Intell.}, Online First.
\end{thebibliography}

\end{document}